% 以下分类谨提供《大连海事大学马克思主义学院博士、硕士学位论文格式规范（20220408）》中的例子，请读者根据实际情况进行编排
\subsubsection*{一、经典著作与文章}
  \bibitem{}{《马克思恩格斯全集》第一至五十卷，北京：人民出版社2016年版。}
  \bibitem{}{《马克思恩格斯选集》第一至四卷，北京：人民出版社2012年版。}
  \bibitem{}{《列宁选集》第一至四卷，北京：人民出版社2012年版。}
  \bibitem{}{《毛泽东选集》第一至四卷，北京：人民出版社1991年版。}
  \bibitem{}{《毛泽东文集》第一至八卷，北京：人民出版社1993年版。}
  \bibitem{}{《毛泽东著作专题摘编》，北京：中央文献出版社2003年版。}
  \bibitem{}{《邓小平文选》第一至三卷，北京：人民出版社1993-1994年版。}
  \bibitem{}{《江泽民文选》第一至三卷，北京：人民出版社2006年版。}
  \bibitem{}{《胡锦涛文选》第一至三卷，北京：人民出版社2016年版。}
  \bibitem{}{习近平：《摆脱贫困》，福州：福建人民出版社2014年版。}
  \bibitem{}{习近平：《干在实处，走在前列：推进浙江新发展的思考与实践》，北京：中共中央党校出版社2006年版。}
  \bibitem{}{习近平：《紧紧围绕坚持和发展中国特色社会主义，学习宣传贯彻党的十八大精神》，《人民日报》2012年11月19日。}
  \bibitem{}{习近平：《关于〈中共中央关于制定国民经济和社会发展第十三个五年规划的建议〉的说明》，《人民日报》2015年11月3日。}
  \bibitem{}{习近平：《在全国政协新年茶话会上的讲话》，《人民日报》2018年12月30日。}
\nextbibtype
\subsubsection*{二、党和政府文献及文献汇编}
  \bibitem{}{《十一届三中全会以来党和国家重要文献选编》，北京：中共中央党校出版社2008年版。}
  \bibitem{}{《十二大以来重要文献选编》（上、中、下），北京：中央文献出版社2011年版。}
  \bibitem{}{《十三大以来重要文献选编》（上、中、下），北京：中央文献出版社2011年版。}
  \bibitem{}{《十四大以来重要文献选编》（上、中、下），北京：中央文献出版社2011年版。}
  \bibitem{}{《十五大以来重要文献选编》（上、中、下），北京：中央文献出版社2011年版。}
  \bibitem{}{《十六大以来重要文献选编》（上、中、下），北京：中央文献出版社2011年版。}
  \bibitem{}{《十七大以来重要文献选编》（上、中、下），北京：中央文献出版社2013年版。}
  \bibitem{}{《十八大以来重要文献选编》（上、中、下），北京：中央文献出版社2014-2018年版。}
\nextbibtype
\subsubsection*{三、中文著作} %按重要程度排序，或者按作者汉语拼音排序，或者按出版时间排序
  \bibitem{}{白暴力：《现代西方经济理论研究》，北京：经济科学出版社2011年版。}
  \bibitem{}{[美]亚当马科斯、马库斯布拉泽顿编，王翰民、杨文锐译：《太平洋之声》，重庆；重庆出版社2014年版。}
\nextbibtype
\subsubsection*{四、中文期刊} %按重要程度排序，或者按作者汉语拼音排序，或者按发表时间排序
  \bibitem{}{张云德、王丽：《习近平对马克思主义中国化理论创新的逻辑起点》，《广西社会科学》2016年第8期。}
\nextbibtype
\subsubsection*{五、外文文献} %按重要程度排序，或者按作者英文首字母排序，或者按出版时间排序
  \bibitem{}{Amrita Jash, China’s “One Belt, One Road”: A Roadmap to“Chinese Dream”?, IndraStra Global, 2016.}
\nextbibtype
\subsubsection*{六、学位论文} %按重要程度排序，或者按作者汉语拼音排序，或者按发表时间排序
  \bibitem{}{李韬：《美国的慈善基金会与美国政治》，中国社会科学院2003年博士学位论文。}
\nextbibtype
\subsubsection*{七、网络文献}
  \bibitem{}{张胜男：《从文化旅游走向创意旅游》，上网日期：2013年5月21日，\url{http://www.wenming.cn/whtzgg_pd/tsyj/201305/t20130521_1238641.shtml}。}
